\section{Introduction}

The rapid democratization of Large Language Models (LLMs) has created a fragmented landscape defined by competing axes of capability, cost, and latency. While frontier models (e.g., GPT-4o, Claude 3.5 Sonnet) offer state-of-the-art reasoning, they are often prohibitively expensive for high-throughput applications. Conversely, efficient models (e.g., Llama-3-8B) are cost-effective but prone to failure on complex tasks. This capability asymmetry necessitates \textbf{LLM Routing}: the dynamic assignment of each query to the most efficient model capable of solving it.

In commercial deployment settings---such as coding assistants or customer support chatbots---this routing decision is constrained by strict Service Level Objectives (SLOs). Operators must often guarantee Time-to-First-Token (TTFT) latency below strict thresholds while keeping monthly inference spend within budget caps, all without degrading task success rates. Existing solutions struggle to satisfy these competing operational constraints simultaneously.

\subsection{The Routing Trilemma}
Current routing architectures impose severe trade-offs that limit their adoption in production environments:

\begin{itemize}
    \item \textbf{Latency Penalties (Cascades):} Approaches like FrugalGPT optimize cost by sequentially querying models until a confidence threshold is met \cite{chen2023frugalgpt}. While effective, this introduces variable latency that grows with the number of fallback calls, often violating real-time SLOs.
    \item \textbf{Supervision Bottlenecks (Predictors):} Supervised routers like RouteLLM rely on massive datasets of human preference labels (e.g., Chatbot Arena) for training \cite{ong2024routellm}. This creates a static policy that adapts to domain shifts (e.g., new proprietary coding standards) only via expensive retraining and fresh labeling cycles.
    \item \textbf{Granularity Blind-Spots (Semantic Maps):} Static embedding routers cluster queries by topic (e.g., ``Math''), assuming uniform difficulty within a cluster. They often fail to distinguish between trivial arithmetic and complex proofs, leading to systematic over-provisioning and measurable budget wastage on evaluation logs.
\end{itemize}

\subsection{Why Contextual Bandits?}
To address these limitations, we formulate routing as a \textbf{Contextual Linear Bandit} problem. Unlike supervised classifiers, bandits naturally handle the partial feedback inherent in production systems: we observe the outcome (success/failure) only for the selected model, never the counterfactuals.

While recent works such as BaRP \cite{ruan2024barp} and PILOT \cite{shin2023pilot} have explored bandit-based routing to optimize scalar performance-cost trade-offs, they typically suffer from the ``Cold Start'' problem---performing near-randomly in the earliest rounds until sufficient data accumulates. Furthermore, many rely on deep reinforcement learning policies with high sample complexity.

In contrast, our focus is on single-shot routing under strict latency constraints with an explicit cold-start mitigation mechanism. We propose a lightweight, interpretable linear framework initialized via \textbf{Generative Calibration}: prior to deployment, the router’s covariance matrix is pre-conditioned using synthetic prompts labeled by a teacher model. This transfers structural knowledge about prompt difficulty (e.g., prompts with formal mathematical notation may systematically differ in difficulty) into the bandit’s prior, allowing it to deploy with near-optimal performance from the very first request while retaining the plasticity to learn from real user feedback.

\subsection{Our Contribution: BanditGPT}
We present \textbf{BanditGPT}, an adaptive routing framework designed for latency-constrained production environments. Our primary contributions are:

\begin{itemize}
    \item \textbf{Elimination of Cold Start} via a procedural warm-up method that uses $\sim$20,000 synthetic samples to pre-condition the bandit's feature geometry, matching the performance of static SOTA baselines (see Section 4.1) within 50 interactions on realistic replay logs.
    \item \textbf{Cost-Efficient Pareto Frontier} achieved through a 37\% reduction in aggregate token costs compared to static best-arm policies on our stratified evaluation set, while maintaining a win-rate within 2\% of the oracle.
    \item \textbf{Single-Shot Latency} with a deterministic routing overhead on the order of 10ms, satisfying real-time latency requirements that cascading architectures often violate.
\end{itemize}

Together, these properties make BanditGPT suitable for real-time, cost-sensitive LLM applications where both budget and latency are tightly constrained.
